\documentclass[10pt,a4paper]{article}
\usepackage[T1]{fontenc}
\usepackage{lmodern}
\usepackage{amsmath,amssymb,amsthm,mathtools}
\usepackage[margin=24mm]{geometry}
\usepackage{microtype}
\usepackage{booktabs,array,longtable}
\usepackage{xcolor}
\usepackage{hyperref}
\usepackage{fancyhdr}
\hypersetup{colorlinks=true,linkcolor=black,citecolor=black,urlcolor=blue,pdfauthor={Matthew J. Colbrook}}
\pagestyle{fancy}
\fancyhf{}
\fancyhead[L]{\small Resolution / MI-23}
\fancyhead[R]{\small 11 September 2026}
\fancyfoot[C]{\thepage}
\setlength{\headheight}{14pt}
\setlength{\parskip}{4pt}
\setlength{\parindent}{0pt}
\newtheorem{theorem}{Theorem}[section]
\newtheorem{proposition}[theorem]{Proposition}
\newtheorem{lemma}[theorem]{Lemma}
\newtheorem{corollary}[theorem]{Corollary}
\theoremstyle{remark}
\newtheorem{remark}[theorem]{Remark}
\newcommand{\M}{\mathbb M}
\newcommand{\C}{\mathbb C}
\newcommand{\R}{\mathbb R}
\newcommand{\tr}{\operatorname{tr}}
\newcommand{\diag}{\operatorname{diag}}
\newcommand{\spec}{\operatorname{spec}}
\newcommand{\rank}{\operatorname{rank}}
\newcommand{\abs}[1]{\lvert #1\rvert}
\newcommand{\norm}[1]{\lVert #1\rVert}
\newcommand{\opnorm}[1]{\lVert #1\rVert_{\infty}}
\newcommand{\schatten}[2]{\lVert #1\rVert_{#2}}
\newcommand{\symmod}{\mathcal S}
\newcommand{\maxmod}{\mathcal M}
\newcommand{\draftnotice}{\begin{center}\small\textit{Proposed mathematical resolution. Complete proof supplied; not submitted or peer reviewed.}\end{center}}

\usepackage{xurl}
\setlength{\emergencystretch}{3em}
\allowdisplaybreaks[1]

\title{MI-23: Resolution}
\author{Matthew J. Colbrook\\
\small Department of Applied Mathematics and Theoretical Physics\\
\small University of Cambridge, Cambridge, United Kingdom\\
\small\href{mailto:m.colbrook@damtp.cam.ac.uk}{m.colbrook@damtp.cam.ac.uk}}
\date{11 September 2026}
\begin{document}
\maketitle
\textbf{Repository status: Solved.}
The complete argument received an independent Codex-agent PASS review
against the exact catalog target.
The original draft was AI-assisted; the reviewed proof below is unchanged.
This is independent agent verification, not external human peer review or
formal proof certification. \href{https://github.com/ajt60gaibb/OpenProblemsInNLA/blob/main/references/colbrook-matrix-2026-09-11/verification/reviews/MI-23-review.md}{Detailed review and scope}.
\par\medskip
% BEGIN REVIEWED PROOF
\section{MI-23: an exact counterexample to the generalized geometric-mean conjecture}
\label{sec:mi23}
For positive definite matrices, write
\[
 A\#_{r,t}B=A^{r/2}(A^{-1/2}BA^{-1/2})^tA^{r/2}.
\]
Conjecture 2.1 of \cite{GAMA2021}, reproduced as MI-23, asserts
\begin{equation}\label{eq:mi23-target}
 \lambda\bigl((A\#_{r,t}B)^p(A\#_{s,1-t}B)^p\bigr)
 \prec_{\log}\lambda\bigl(A^{p(r+s-1)}B^p\bigr)
\end{equation}
for $0\leq t\leq1$, $p\geq1$, and either $r,s\geq1$ or $r,s\leq0$. Products of two positive definite matrices have positive eigenvalues, which are decreasingly ordered here. The assertion fails already for $r=s=1$ and $p=2$.

\begin{theorem}\label{thm:mi23}
Let
\[
 D=\diag\left(16,\frac1{12},1\right),\qquad
 T=\begin{pmatrix}2&1&2\\1&25&-10\\2&-10&10\end{pmatrix},
 \qquad A=D^2,\qquad B=DT^8D.
\]
Then $A,B>0$, and \eqref{eq:mi23-target} is false at
\[
 r=s=1,\qquad p=2,\qquad t=\frac18.
\]
\end{theorem}
\begin{proof}
The leading principal minors of $T$ are $2$, $49$, and $150$. Thus $T>0$, and consequently $A,B>0$. Moreover,
\[
 A^{-1/2}BA^{-1/2}=T^8.
\]
Its positive eighth root is exactly $T$, not merely a numerical approximation. Hence the two geometric means are the rational matrices
\[
 G=A\#_{1,1/8}B=DTD,\qquad
 H=A\#_{1,7/8}B=DT^7D.
\]
The first inequality in the proposed log-majorization would imply
\[
 \lambda_{\max}(G^2H^2)\leq\lambda_{\max}(A^2B^2).
\]
Since $G,H,A,B$ are positive definite,
\[
 \lambda_{\max}(G^2H^2)=\opnorm{GH}^2,
 \qquad \lambda_{\max}(A^2B^2)=\opnorm{AB}^2.
\]
For example, $G^2H^2$ is similar to $HG^2H=(GH)^*(GH)$, which proves the first identity; the second is identical.

Exact rational matrix multiplication gives
\begin{align*}
 (GH)_{13}&=\frac{1260589125202}{9}>140000000000,\\
 \norm{AB}_F^2&=\frac{2009446159144992718181231562721}{107495424}
 <(138000000000)^2.
\end{align*}
It follows that
\[
 \opnorm{GH}\geq |(GH)_{13}|>140000000000
 >138000000000>\norm{AB}_F\geq\opnorm{AB},
\]
contradicting the required largest-eigenvalue inequality.
For an alternative one-line rational certificate of the strict comparison,
\[
 ((GH)_{13})^2-\norm{AB}_F^2
 =\frac{99434824489435745411095588895}{107495424}>0.
\]
\end{proof}

\begin{remark}[Exactness and parameter scope]
All input matrices and every displayed certificate are rational. The representation $B=DT^8D$ eliminates all fractional-power computations from the proof. This counterexample is outside the narrower interval $1/4\leq t\leq3/4$ established in \cite{GAMA2021} for $p=2$; it does not contradict that theorem. It disproves the stated conjecture for the full interval $0\leq t\leq1$.
\end{remark}

% END REVIEWED PROOF
\begin{thebibliography}{1}

\bibitem{GAMA2021}
Mohammad~M. Ghabries, Hassane Abbas, Bassam Mourad, and Abdallah Assi.
\newblock New log-majorization results concerning eigenvalues and singular
  values and a complement of a norm inequality.
\newblock arXiv:2105.13356, 2021.
\newblock Conjectures 1.1 and 2.1.
\newblock URL: \url{https://arxiv.org/abs/2105.13356}.

\end{thebibliography}

\end{document}
